% Version of June 15, 2021
%
%%%%%%%%%%%%%%%%%%%%%%%%%%%%%%%%%%%%%%%%%%%%%%%%%%%%%%%%%%%%%%%%%%%%%%%%%%%%%
% The first set of commands control the overall configuration - such as,
% is color used to show changes since the last version, are changebars
% used, etc.
% Control the use of color
\newif\ifusecolor
\usecolortrue
% For publisher's additions in the printed book
\newif\ifbookprinting
\bookprintingfalse
% Provide a way to indicate whether the text is within a float (such as a
% table); this is necessary to avoid problems with lost floats.  Use this
% for macros that use the \ticket command (such as the MPIupdate, MPIreplace,
% or MPIdelete macros) within table for figure environments.
\newif\ifinfloat
\infloatfalse

%
% Use this to make pages start on the right side
\newcommand{\startchap}[0]{\cleardoublepage}
%\newcommand{\startchap}[0]{\relax}

%%%%%%%%%%%%%%%%%%%%%%%%%%%%%%%%%%%%%%%%%%%%%%%%%%%%%%%%%%%%%%%%%%%%%%%%%%%%%
% There are two elements to ``draft'':
%   1. Mark the headers/footers/watermark with draft information
%   2. Provide text about the ``draft specification''
% use \drafttrue for the headers/footers/watermaky
% use \draftspectrue to turn on the text abot the draft specification
\newif\ifdraftspec
\draftspecfalse

\newif\ifdraft
\def\drafttext{\upshape\normalsize\textbf{Unofficial Draft for Comment Only}}
\def\insertdraft{\ifdraft\drafttext\else\fi}


%%%%%%%%%%%%%%%%%%%%%%%%%%%%%%%%%%%%%%%%%%%%%%%%%%%%%%%%%%%%%%%%%%%%%%%%%%%%%
% The following places line numbers along the page sides
%
% Get line numbers in the gutters.  Thanks to Guy Steele and HPFF!
%

\makeatletter
%
% This is used to put line numbers on plain pages.  Used in draft.tex
%
\def\withlinenumbers{\relax
  \def\@evenfoot{\hbox to 0pt{\hss\LineNumberRuler\hskip 1.5pc}\hfil\insertdraft\hfil}\relax
  \def\@oddfoot{\hfil\insertdraft\hfil\hbox to 0pt{\hskip 1.5pc\LineNumberRuler\hss}}}

\def\LineNumberRuler{\vbox to 0pt{\vss\normalsize \baselineskip13.6pt
    \lineskip 1pt \normallineskip 1pt \def\baselinestretch{1}\relax
    \LNR{1}\LNR{2}\LNR{3}\LNR{4}\LNR{5}\LNR{6}\LNR{7}\LNR{8}\LNR{9}
    \LNR{10}\LNR{11}\LNR{12}\LNR{13}\LNR{14}
        \LNR{15}\LNR{16}\LNR{17}\LNR{18}\LNR{19}
    \LNR{20}\LNR{21}\LNR{22}\LNR{23}\LNR{24}
        \LNR{25}\LNR{26}\LNR{27}\LNR{28}\LNR{29}
    \LNR{30}\LNR{31}\LNR{32}\LNR{33}\LNR{34}\LNR{35}
        \LNR{36}\LNR{37}\LNR{38}\LNR{39}
    \LNR{40}\LNR{41}\LNR{42}\LNR{43}\LNR{44}
        \LNR{45}\LNR{46}\LNR{47}\LNR{48}
    \vskip 31pt
%    \vskip 20pt % Additional adjustment when using scrbook/latex-headings
    \vskip 17pt % Additional ajustment when using srcbook/headings
}}
\def\LNR#1{\hbox to 1pc{\hfil\tiny#1\hfil}}

% jmm; merge the withlinenumbers stuff into
% the centered page numbers that tex defines by default
\def\ps@plainwithlinenumbers{\let\@mkboth\@gobbletwo
     \def\@oddhead{}
     \def\@oddfoot{\hfil\rmfamily\ifdraft\drafttext\hfil\thepage\else\thepage\hfil\fi
       \hbox to 0pt{\hskip 1.5pc\LineNumberRuler\hss}}
     \def\@evenhead{}
     \def\@evenfoot{\hbox to 0pt{\hss
     \LineNumberRuler\hskip 1.5pc}\rmfamily\hfil\ifdraft\drafttext\hfil\thepage\else\thepage\hfil\fi}}

% The old version; uncommenting the withlinenumbers part didn't
% work because that macro replaced the footer definition that
% did page numbering.
%\def\ps@plainwithlinenumbers{\ps@plain}%\withlinenumbers}
% end jmm changes
%%%%%%%%%%%%%%%%%%%%%%%%%%%%%%%%%%%%%%%%%%%%%%%%%%%%%%%%%%%%%%%%%%%%%%%%%%%%%

%
% 1st page of a chapter has its own page style, so we have to put line
% numbers in here also.
%
\newwrite\chappages
\immediate\openout\chappages=chappage.txt
\def\writespace{ }

%%%%%%%%%%%%%%%%%%%%%%%%%%%%%%%%%%%%%%%%%%%%%%%%%%%%%%%%%%%%%%%%%%%%%%%%%%%%%
%
% Contents is done with \chapter*{Contents}, so we need to turn off the
% line numbers in this case.  Easiest to look at def
%
\def\incontents{0}
\newif\ifcontents
\contentsfalse
\def\chapter{\clearpage \ifcontents\else\thispagestyle{plainwithlinenumbers}\fi
        \write\chappages{Chapter \thechapter\writespace - \the\count0}
        \global\@topnum\z@ \@afterindentfalse \secdef\@chapter\@schapter}

%
% Change "Chapter" to "Chapter", "Appendix" to "Annex"
%
\renewcommand{\chaptername}{Chapter}
\renewcommand{\appendixname}{Annex}
% ... old code does not work correctly with pdflatex
% \def\@chapapp{Chapter}
% \def\appendix{\par
%  \setcounter{chapter}{0}
%  \setcounter{section}{0}
%  \def\@chapapp{Annex}
%  \def\thechapter{\Alph{chapter}}}

\makeatother
%
% Use Sans Serif font for sections, etc.  S. Otto
%
\if 0 % Use standard section etc as defined by book or srcbook style
\makeatletter
\def\section{\@startsection {section}{1}{\z@}{-3.5ex plus -1ex minus 
-.2ex}{2.3ex plus .2ex}{\Large\sf}}
\def\subsection{\@startsection{subsection}{2}{\z@}{-3.25ex plus -1ex minus 
-.2ex}{1.5ex plus .2ex}{\large\sf}}
\def\subsubsection{\@startsection{subsubsection}{3}{\z@}{-3.25ex plus 
-1ex minus -.2ex}{1.5ex plus .2ex}{\normalsize\sf}}
\def\paragraph{\@startsection {paragraph}{4}{\z@}{3.25ex plus 1ex 
minus .2ex}{-1em}{\normalsize\sf}}
\makeatother
\fi

%%%%%%%%%%%%%%%%%%%%%%%%%%%%%%%%%%%%%%%%%%%%%%%%%%%%%%%%%%%%%%%%%%%%%%%%%%%%%
%%
%% For the funcdef macros, we use s list environment.  To make this consistent,
%%
\def\MPIfunclist{
    \begin{list}{}{                     % see pg 113 of Lamport's book
        \setlength{\leftmargin}{200pt}
        \setlength{\labelwidth}{180pt}
        \setlength{\labelsep}{10pt}
        \setlength{\itemindent}{0pt}
        \setlength{\itemsep}{0pt}
        \setlength{\topsep}{5pt}
    }
}

%%%%%%%%%%%%%%%%%%%%%%%%%%%%%%%%%%%%%%%%%%%%%%%%%%%%%%%%%%%%%%%%%%%%%%%%%%%%%
\makeatletter
%
% make our own index environment that can have a different
% title than just "Index" -- S. Otto
%
\def\@index{Index}
\def\@introtext{}  % MPI-2.1
\newif\if@restonecol
\def\myindex{\@restonecoltrue\if@twocolumn\@restonecolfalse\fi
\columnseprule \z@
%\columnsep 35pt\twocolumn[\@makeschapterhead{Index}]
 %\@mkboth{INDEX}{INDEX}\thispagestyle{plain}\parindent\z@
%\columnsep 35pt\twocolumn[\@makeschapterhead{\@index}]
%\columnsep 35pt\twocolumn[\@makeschapterhead{\@index}\@introtext\vspace{15pt}] % MPI-2.1
\columnsep 20pt\twocolumn[\@makeschapterhead{\@index}\@introtext\vspace{15pt}] % MPI-3.0
 \@mkboth{\@index}{\@index}\thispagestyle{plainwithlinenumbers}\parindent\z@
 \parskip\z@ plus .3pt\relax\let\item\@idxitem}
\def\@idxitem{\par\hangindent 40pt}
\def\subitem{\par\hangindent 40pt \hspace*{20pt}}
\def\subsubitem{\par\hangindent 40pt \hspace*{30pt}}
\def\endmyindex{\if@restonecol\onecolumn\else\clearpage\fi}
\def\indexspace{\par \vskip 10pt plus 5pt minus 3pt\relax}
\makeatother

\newcommand{\uu}[1]{\underline{\hyperpage{#1}}} 
\newcommand{\bold}[1]{\textbf{\hyperpage{#1}}} 
%%%%%%%%%%%%%%%%%%%%%%%%%%%%%%%%%%%%%%%%%%%%%%%%%%%%%%%%%%%%%%%%%%%%%%%%%%%%%


%%%%%%%%%%%%%%%%%%%%%%%%%%%%%%%%%%%%%%%%%%%%%%%%%%%%%%%%%%%%%%%%%%%%%%%%%%%%%
% These improve the formatting of backref links in the bibliography
% See http://ctan.math.washington.edu/tex-archive/macros/latex/contrib/hyperref/doc/manual.pdf , section 5.26 (New Featurs/Backref entries)
% The regular backref features don't work when combined with the hyperref
% package
\def\fixupbackrefformatting{%
\renewcommand*{\backref}[1]{
% default interface
% #1: backref list
%
% We want to use the alternative interface,
% therefore the definition is empty here.
}
\renewcommand*{\backrefalt}[4]{
\ifnum ##1=1 %
 Citation on page ##2.
\else
 Citations on pages ##2.
\fi}
}
%%%%%%%%%%%%%%%%%%%%%%%%%%%%%%%%%%%%%%%%%%%%%%%%%%%%%%%%%%%%%%%%%%%%%%%%%%%%%

\raggedbottom

% deleting the content where the macro is used; but preserving it for the Annex
\def\OnlyForAutomaticAnnexGeneration#1{}
